\subsection{Tree extensions of categorical acts}

\begin{lemma}[Proved: compatible trees]
\label{A7:tree-act}
Let $F$ be a nonempty finite-set-valued left or right act on a
finite category $S$.  Suppose that, on each set $F(s)$, one can
choose a tree $T(s)$ with exactly $F(s)$ as its vertex set, so
that every act map sends the endpoints of each tree edge to
equal or adjacent vertices.  Then every positive rational
cohomology product of the act-cone category is zero.
\end{lemma}
\begin{proof}
Give every edge length one. Extend each act map linearly on an
edge whose endpoints have distinct images, and constantly on
an edge whose endpoints have the same image. The hypothesis
makes these continuous maps of trees. They compose strictly:
an edge is sent either to a vertex or linearly to a single
edge, and its subsequent image is again determined by its
endpoints. Thus the extensions form an actual diagram $T$,
with a natural vertex inclusion $F\hookrightarrow T$.

Every $T(s)$ is contractible. Hence the natural map
$\operatorname{hocolim}T\to BS$ is a homotopy equivalence:
on each level of the simplicial bar it is a disjoint union of
products of a contractible tree with a simplex. The act cone
is therefore equivalent to the homotopy cofiber of
$\operatorname{hocolim}F\to\operatorname{hocolim}T$.
The objectwise inclusions are cellular cofibrations. This
cofiber is the pointed bar of the quotients $T(s)/F(s)$, with
the bar of their basepoints collapsed, exactly as in
Lemma~\ref{CP2:A-high-affine}.

Each $T(s)/F(s)$ is a wedge of one circle for each edge of
$T(s)$. Each actor sends any such circle to the basepoint or
linearly to one target circle, with degree $1$ or $-1$.
Choose an orientation on every tree edge. The degree-two maps
on all these circles commute strictly with every actor,
including orientation reversal, since $2(-t)=-2t$ modulo one.
They therefore induce a self-map of the pointed bar.
The reduced cellular chain complex of every quotient has
only degree one, and the self-map is multiplication by two
there. It is consequently multiplication by two on the
entire reduced bar total chain complex and on every positive
cohomology group. For positive classes $u,v$, naturality
gives $2uv=4uv$ over $\mathbb Q$, so $uv=0$.
The same construction with the opposite variance proves
both left- and right-act statements.
\end{proof}

\begin{proposition}[Proved: small trees at the source or sink]
\label{A7:tree-stalk-profile}
Let $\mathcal C$ be a finite acyclic category with full support
$0<1<2<3<4$, and put $h_{ij}=|\mathcal C(i,j)|$.
Every rational degree-two cohomology product of $B\mathcal C$
is zero if either of the following source conditions holds:
\[
\begin{split}
 &h_{02}\leq2,\qquad h_{03}\leq2;\qquad\text{(S1)}\\
 &h_{01}\leq2,\quad h_{03}\leq2,
       \quad h_{12}\leq3.\qquad\text{(S2)}
\end{split}
\]
The same conclusion holds under either dual sink condition
\[
\begin{split}
 &h_{24}\leq2,\qquad h_{14}\leq2;\qquad\text{(T1)}\\
 &h_{34}\leq2,\quad h_{14}\leq2,
       \quad h_{23}\leq3.\qquad\text{(T2)}
\end{split}
\]
There is no restriction on $h_{04}$.
\end{proposition}
\begin{proof}
Collapse the extreme hom-set $\mathcal C(0,4)$ to a singleton.
Proposition~\ref{A7:extreme-arrow-collapse} gives a surjection
on second cohomology from this quotient category to the
original category, preserving products. It suffices to work
in the quotient, which is the left-act cone of its full
subcategory on $1,2,3,4$. Write its source stalks as
$A=\mathcal C(0,1)$, $U=\mathcal C(0,2)$,
$L=\mathcal C(0,3)$ and the singleton at $4$.

Under (S1), choose any tree on $A$. Give each of $U,L$ its
unique tree with at most two vertices, and use the point at
$4$. Every act map carries an edge to an edge or a vertex:
all maps from $A$ have target of size at most two, as do the
remaining nonidentity maps. Lemma~\ref{A7:tree-act} proves
the conclusion, with no bound on $|A|$.

For (S2), write $B=\mathcal C(1,2)$. On the full vertex set
$U$, form a simple undirected graph $G$ by adding, for
each $b\in B$, the pair of distinct values of the map
$A\xrightarrow{\,b\circ-\,}U$, if that map has two values.
There are at most three edges. If $G$ is a forest, extend it
to a tree on $U$. Use the unique tree on $A$ and the unique
tree on $L$, both of which have at most two vertices.
Every map $A\to U$ sends the only possible edge to an edge
of the chosen tree or to a vertex. Every other act map has
target $L$ or the singleton at $4$, so it automatically has
the same property. Lemma~\ref{A7:tree-act} again applies.

The only remaining possibility is that $G$ is a triangle
together with isolated vertices: a simple graph with at
most three edges has no other cycle.
Then $|A|=2$, $|B|=3$, and the three maps $A\to U$ give its
three distinct edges. Let $A=\{a_0,a_1\}$ and write
$P(a,b)=b\circ a$. Any rational two-cycle on the full
subcategory with objects $0,1,2$ has, from its $B$-edge
boundary coefficients, the form
\[
 z=\sum_{b\in B}t_b\bigl[(a_0,b)-(a_1,b)\bigr].
\]
Its remaining boundary equations include
\[
 \sum_b t_b=0,
 \qquad
 \sum_b t_b\bigl([P(a_1,b)]-[P(a_0,b)]\bigr)=0.
\]
The second equation is the incidence-kernel equation of
the triangle with its three independently chosen edge
orientations; any isolated vertices of $U$ contribute
only zero rows. Its kernel is spanned by a vector whose
three entries are all $1$ or $-1$. Their sum is an odd
integer and cannot be zero over $\mathbb Q$. The first
equation therefore forces every $t_b=0$. The full
subcategory on $012$ has $H_2=0$, and hence $H^2=0$,
over $\mathbb Q$.

For any two global degree-two cocycles $\alpha,\beta$,
the restriction of $\alpha$ to $012$ is consequently a
coboundary. Extend a primitive from the edges of $012$
by zero and subtract its global coboundary. The resulting
representative of $\alpha$ vanishes on every $012$
two-simplex. Its Alexander--Whitney product with $\beta$
is zero on every four-simplex, since every such simplex
uses the object string $01234$. This proves (S2).
Opposition gives (T1) and (T2).
\end{proof}

\begin{remark*}
The restriction $h_{12}\leq3$ in (S2) is used only in
the triangle case. With more arrows in that hom-set,
repeated oriented edges can give incidence-kernel vectors
with coefficient sum zero, and the displayed argument
does not assert the same cardinality reduction.
\end{remark*}
